\documentclass[../../main.tex]{subfiles}
\begin{document}
\section{光滑覆叠映射}\label{section:ch04-光滑覆叠映射}

本节我们引进一类在光滑流形理论中扮演重要角色的局部微分同胚。你大概已经熟悉拓扑空间之间覆叠映射的概念:它是连通、局部道路连通空间之间的满射连续映射 $\pi:E\to M$,使得 $M$ 的每一点都有一个被\textbf{均匀覆盖}的邻域 $U$,即 $\pi^{-1}(U)$ 的每个分支都被 $\pi$ 同胚地映到 $U$ 上。覆叠映射的基本性质汇总于附录 A(见原书第 615--616 页)。

在光滑流形的背景下,引进一种稍受限的覆叠映射是有用的。若 $E$ 与 $M$ 是带边或不带边的连通光滑流形,称映射 $\pi:E\to M$ 为\textbf{光滑覆叠映射},如果 $\pi$ 光滑且满射,并且 $M$ 中每一点都有邻域 $U$,使 $\pi^{-1}(U)$ 的每个分支都被 $\pi$ 微分同胚地映到 $U$ 上。在这一背景下我们也称 $U$ 被均匀覆盖。空间 $M$ 称为覆叠的底空间, $E$ 称为 $M$ 的覆叠流形。若 $E$ 单连通,则它称为 $M$ 的万有覆叠流形。

为把这一新定义与前面的定义区分开,我们常把通常的(未必光滑的)覆叠映射称为拓扑覆叠映射。光滑覆叠映射特别地是拓扑覆叠映射。但与本章研究过的其他类型映射一样,光滑覆叠映射不仅仅是恰好光滑的拓扑覆叠映射:定义还要求 $\pi$ 在被均匀覆盖集合的原像的每个分支上的限制是微分同胚,而不仅仅是光滑同胚。

\begin{proposition}[Properties of Smooth Coverings]\label{proposition:ch04-covering-properties}
\begin{enumerate}[(a)]
\item 每个光滑覆叠映射都是局部微分同胚、光滑淹没、开映射与商映射。
\item 单射光滑覆叠映射是微分同胚。
\item 拓扑覆叠映射是光滑覆叠映射,当且仅当它是局部微分同胚。
\end{enumerate}
\end{proposition}

\begin{exercise}\label{exercise:ch04-prove-covering-props}
证明\cref{proposition:ch04-covering-properties}。
\end{exercise}

\begin{example}[Smooth Covering Maps]\label{example:ch04-smooth-coverings}
\cref{example:ch01-circle-covering}(b) 中定义的映射 $\varepsilon:\mathbb{R}\to S^1$ 是拓扑覆叠映射且是局部微分同胚(另见\cref{example:ch04-local-diffeomorphisms}),故它是光滑覆叠映射。类似地, \cref{example:ch01-circle-covering}(c) 中的映射 $\varepsilon^n:\mathbb{R}^n\to\mathbb{T}^n$ 是光滑覆叠映射。对每个 $n\geqslant 1$, \cref{example:ch01-circle-covering}(f) 中定义的映射 $q:S^n\to\mathbb{RP}^n$ 是二叶光滑覆叠映射(见习题 4-10)。
\end{example}

由于光滑覆叠映射是满射光滑淹没,上一节关于光滑淹没的所有结果都可应用于它们。例如, \cref{theorem:ch04-pass-quotient} 是一个特别有用的工具,用于从覆叠空间的底空间出发定义光滑映射。这种技巧的例子见习题 4-12 与习题 4-13。

对光滑覆叠映射,局部截面定理可以加强。

\begin{proposition}[Local Section Theorem for Smooth Covering Maps]\label{proposition:ch04-local-section-covering}
设 $E$ 与 $M$ 为带边或不带边的光滑流形, $\pi:E\to M$ 是光滑覆叠映射。给定任一被均匀覆盖的开子集 $U\subseteq M$、任一 $q\in U$、以及 $\pi$ 在 $q$ 上方纤维中的任一点 $p$,存在唯一光滑局部截面 $\sigma:U\to E$ 使 $\sigma(q)=p$。
\end{proposition}

\begin{proof}
设 $U\subseteq M$ 被均匀覆盖, $q\in U$, $p\in\pi^{-1}(q)$。令 $\widetilde{U}_0$ 为 $\pi^{-1}(U)$ 中含 $p$ 的分支。由于 $\pi$ 在 $\widetilde{U}_0$ 上的限制是到 $U$ 上的微分同胚,映射 $\sigma=\bigl(\pi|_{\widetilde{U}_0}\bigr)^{-1}$ 即为所求的光滑局部截面。

为证唯一性,设 $\sigma':U\to E$ 是任一其他满足 $\sigma'(q)=p$ 的光滑局部截面。由于 $U$ 连通, $\sigma'(U)$ 包含在含 $p$ 的分支 $\widetilde{U}_0$ 中。由于 $\sigma'$ 是双射映射 $\pi|_{\widetilde{U}_0}$ 的右逆,它必等于其逆,因而等于 $\sigma$。
\end{proof}

\begin{exercise}\label{exercise:ch04-covering-section-smooth}
设 $\pi:E\to M$ 是光滑覆叠映射。证明 $\pi$ 的每个局部截面都光滑。
\end{exercise}

\begin{exercise}\label{exercise:ch04-product-coverings}
设 $E_1,\dots,E_k$ 与 $M_1,\dots,M_k$ 为(不带边的)光滑流形,且对每个 $i=1,\dots,k$, $\pi_i:E_i\to M_i$ 是光滑覆叠映射。证明 $\pi_1\times\cdots\times\pi_k:E_1\times\cdots\times E_k\to M_1\times\cdots\times M_k$ 是光滑覆叠映射。
\end{exercise}

\begin{exercise}\label{exercise:ch04-evenly-covered-equiv}
设 $\pi:E\to M$ 是光滑覆叠映射。由于 $\pi$ 也是拓扑覆叠映射,关于子集 $U\subseteq M$ 被均匀覆盖的含义存在一个潜在歧义: $\pi$ 是把 $\pi^{-1}(U)$ 的各分支微分同胚地映到 $U$ 上,还是仅仅同胚地映到 $U$ 上?证明这两个概念等价:若 $U\subseteq M$ 在拓扑意义下被均匀覆盖,则 $\pi$ 把 $\pi^{-1}(U)$ 的每个分支微分同胚地映到 $U$ 上。
\end{exercise}

\begin{proposition}[Covering Spaces of Smooth Manifolds]\label{proposition:ch04-covering-smooth}
设 $M$ 是连通光滑 $n$ 流形, $\pi:E\to M$ 是拓扑覆叠映射。则 $E$ 是拓扑 $n$ 流形,并具有唯一的光滑结构使 $\pi$ 成为光滑覆叠映射。
\end{proposition}

\begin{proof}
由于 $\pi$ 是局部同胚, $E$ 局部 Euclid。为证它 Hausdorff,设 $p_1,p_2$ 为 $E$ 中不同两点。若 $\pi(p_1)=\pi(p_2)$ 且 $U\subseteq M$ 是含 $\pi(p_1)$ 的被均匀覆盖开子集,则 $\pi^{-1}(U)$ 中含 $p_1$ 与 $p_2$ 的分支是 $E$ 中分离 $p_1$ 与 $p_2$ 的不交开子集。另一方面,若 $\pi(p_1)\neq\pi(p_2)$,则存在分别含 $\pi(p_1)$ 与 $\pi(p_2)$ 的不交开子集 $U_1,U_2\subseteq M$,于是 $\pi^{-1}(U_1)$ 与 $\pi^{-1}(U_2)$ 是 $E$ 中含 $p_1$ 与 $p_2$ 的不交开子集。故 $E$ Hausdorff。

为证 $E$ 第二可数,我们先证 $\pi$ 的每条纤维可数。给定 $q\in M$ 与 $\pi^{-1}(q)$ 中任一点 $p_0$,我们构造满射 $\beta:\pi_1(M,q)\to\pi^{-1}(q)$;由\cref{proposition:ch01-fundamental-group-countable} 知 $\pi_1(M,q)$ 可数,这就足够了。令 $[f]\in\pi_1(M,q)$ 为基于 $q$ 的任一环路 $f:[0,1]\to M$ 的路径类。覆叠映射的路径提升性质(\cref{proposition:appA-path-lifting})保证存在 $f$ 的提升 $\widetilde{f}:[0,1]\to E$ 从 $p_0$ 出发,而单值性定理(\cref{proposition:appA-monodromy})表明端点 $\widetilde{f}(1)\in\pi^{-1}(q)$ 只依赖于 $f$ 的路径类,故可定义 $\beta[f]=\widetilde{f}(1)$。为看 $\beta$ 满射,只需注意对任一点 $p\in\pi^{-1}(q)$,存在 $E$ 中从 $p_0$ 到 $p$ 的路径 $\widetilde{f}$,于是 $f=\pi\circ\widetilde{f}$ 是 $M$ 中的环路使 $p=\beta[f]$。

所有被均匀覆盖开子集之集是 $M$ 的开覆盖,因而有可数子覆盖 $\{U_i\}$。对任一给定的 $i$, $\pi^{-1}(U_i)$ 的每个分支在 $U_i$ 上方的每条纤维中恰含一点,故 $\pi^{-1}(U_i)$ 有可数多个分支。于是形如 $\pi^{-1}(U_i)$ 的所有集合的所有分支之集是 $E$ 的可数开覆盖;由于每个这样的分支第二可数,由\cref{exercise:appA-second-countable} 知 $E$ 第二可数。这就完成了 $E$ 是拓扑流形的证明。

为在 $E$ 上构造光滑结构,设 $p$ 为 $E$ 中任一点,并设 $U$ 是 $\pi(p)$ 的一个被均匀覆盖邻域。必要时缩小 $U$,可设它也是某光滑坐标映射 $\varphi:U\to\mathbb{R}^n$ 的定义域(\cref{figure:ch04-covering-charts})。若 $\widetilde{U}$ 是 $\pi^{-1}(U)$ 中含 $p$ 的分支,且 $\widetilde{\varphi}=\varphi\circ\pi|_{\widetilde{U}}:\widetilde{U}\to\mathbb{R}^n$,则 $(\widetilde{U},\widetilde{\varphi})$ 是 $E$ 上的坐标卡。若两个这样的坐标卡 $(\widetilde{U},\widetilde{\varphi})$ 与 $(\widetilde{V},\widetilde{\psi})$ 相交,则转移映射可写为
\begin{align*}
\widetilde{\psi}\circ\widetilde{\varphi}^{-1}
&=\bigl(\psi\circ\pi|_{\widetilde{U}\cap\widetilde{V}}\bigr)\circ\bigl(\varphi\circ\pi|_{\widetilde{U}\cap\widetilde{V}}\bigr)^{-1}\\
&=\psi\circ\bigl(\pi|_{\widetilde{U}\cap\widetilde{V}}\bigr)\circ\bigl(\pi|_{\widetilde{U}\cap\widetilde{V}}\bigr)^{-1}\circ\varphi^{-1}\\
&=\psi\circ\varphi^{-1},
\end{align*}
它光滑。于是所有这样的坐标卡之集在 $E$ 上定义了一个光滑结构。这个光滑结构的唯一性留给读者(习题 4-9)。最后, $\pi$ 是光滑覆叠映射,因为用上述构造的任一对坐标卡 $(\widetilde{U},\widetilde{\varphi})$ 与 $(U,\varphi)$ 表示时它的坐标表示是恒等映射。
\end{proof}

\begin{figure}[ht]
\centering
\includegraphics[width=0.8\textwidth]{fig-04-5.png}
\caption{覆盖流形上图的光滑相容性}\label{figure:ch04-covering-charts}
\end{figure}

下面是带边流形的相应结果。

\begin{proposition}[Covering Spaces of Smooth Manifolds with Boundary]\label{proposition:ch04-covering-boundary}
设 $M$ 是连通光滑 $n$ 流形带边, $\pi:E\to M$ 是拓扑覆叠映射。则 $E$ 是拓扑 $n$ 流形带边,满足 $\partial E=\pi^{-1}(\partial M)$,并具有唯一的光滑结构使 $\pi$ 成为光滑覆叠映射。
\end{proposition}

\begin{exercise}\label{exercise:ch04-prove-covering-boundary}
证明上述命题。
\end{exercise}

\begin{corollary}[Existence of a Universal Covering Manifold]\label{corollary:ch04-universal-covering}
若 $M$ 是连通光滑流形,则存在单连通光滑流形 $\widetilde{M}$(称为 $M$ 的万有覆叠流形)与光滑覆叠映射 $\pi:\widetilde{M}\to M$。万有覆叠流形在如下意义下唯一:若 $\widetilde{M}'$ 是任一其他容许光滑覆叠映射 $\pi':\widetilde{M}'\to M$ 的单连通光滑流形,则存在微分同胚 $\Phi:\widetilde{M}\to\widetilde{M}'$ 使 $\pi'\circ\Phi=\pi$。
\end{corollary}

\begin{exercise}\label{exercise:ch04-prove-universal-covering}
证明上述推论。
\end{exercise}

\begin{exercise}\label{exercise:ch04-universal-covering-boundary}
把上述推论推广到带边光滑流形。
\end{exercise}

判断一个给定映射是否为光滑覆叠映射的简单判别法并不多,即使已知它是满射局部微分同胚。下面的命题给出一个有用的充分判别法。(然而它不是必要的;见习题 4-11。)

\begin{proposition}\label{proposition:ch04-proper-covering}
设 $E$ 与 $M$ 是非空连通带边或不带边光滑流形。若 $\pi:E\to M$ 是固有局部微分同胚,则 $\pi$ 是光滑覆叠映射。
\end{proposition}

\begin{proof}
由于 $\pi$ 是局部微分同胚,它是开映射;又由于它固有,它是闭映射(\cref{theorem:appA-proper-closed})。故 $\pi(E)$ 在 $M$ 中既开又闭。由于它显然非空,它就是整个 $M$;故 $\pi$ 满射。

令 $q\in M$ 任意。由于 $\pi$ 是局部微分同胚, $\pi^{-1}(q)$ 中每一点都有邻域使 $\pi$ 在其上单射,故 $\pi^{-1}(q)$ 是 $E$ 的离散子集。由于 $\pi$ 固有, $\pi^{-1}(q)$ 也紧,故它有限。记
\begin{align*}
\pi^{-1}(q)=\{p_1,\dots,p_k\}.
\end{align*}
对每个 $i$,存在 $p_i$ 的邻域 $V_i$,使 $\pi$ 在其上是到某开子集 $U_i\subseteq M$ 上的微分同胚。必要时缩小每个 $V_i$,可设 $V_i\cap V_j=\varnothing$ 当 $i\neq j$。令 $U=U_1\cap\cdots\cap U_k$(\cref{figure:ch04-proper-covering}),它是 $q$ 的邻域。则 $U$ 满足
\begin{align}\label{align:ch04-U-in-Ui}
U\subseteq U_i\quad\text{对每个 }i.
\end{align}
由于 $K=E\setminus(V_1\cup\cdots\cup V_k)$ 在 $E$ 中闭且 $\pi$ 是闭映射, $\pi(K)$ 在 $M$ 中闭。把 $U$ 换成 $U\setminus\pi(K)$,可设 $U$ 还满足
\begin{align}\label{align:ch04-pi-1-U}
\pi^{-1}(U)\subseteq V_1\cup\cdots\cup V_k.
\end{align}
最后,把 $U$ 换成 $U$ 中含 $q$ 的连通分支后,可设 $U$ 连通且仍满足\eqref{align:ch04-U-in-Ui} 与\eqref{align:ch04-pi-1-U}。我们将证明 $U$ 被均匀覆盖。

令 $\widetilde{V}_i=\pi^{-1}(U)\cap V_i$。由\eqref{align:ch04-pi-1-U}, $\pi^{-1}(U)=\widetilde{V}_1\cup\cdots\cup\widetilde{V}_k$。由于 $\pi:V_i\to U_i$ 是微分同胚, \eqref{align:ch04-U-in-Ui} 蕴涵 $\pi:\widetilde{V}_i\to U$ 仍是微分同胚,特别地 $\widetilde{V}_i$ 连通。由于 $\widetilde{V}_1,\dots,\widetilde{V}_k$ 是 $\pi^{-1}(U)$ 的不交连通开子集,它们恰好就是 $\pi^{-1}(U)$ 的各分支。
\end{proof}

\begin{figure}[ht]
\centering
\includegraphics[width=0.8\textwidth]{fig-04-6.png}
\caption{逆紧局部微分同胚是覆盖映射}\label{figure:ch04-proper-covering}
\end{figure}

\end{document}